\section{Reproducibility and evidence ledger}
\label{app:reproducibility-ledger}

The code, configurations, exact audit, and figure scripts for this paper are in the public repository \url{https://github.com/ioannist/six-birds-layer-birth}. This appendix describes how the claims are tied to reproducible outputs.

\paragraph{Evidence ledger.} A reproducibility freeze, written to \path{results/freeze/reproducibility_freeze/}, maps each claim to the result bundles that support it. It reruns every producing workflow with result caches disabled and compares output hashes. Its key files are
\begin{quote}
\path{analysis/evidence_ledger.json}\\
\path{analysis/claim_coverage.json}\\
\path{analysis/reproducibility_summary.json}\\
\path{analysis/ready_for_writing_status.json}.
\end{quote}
Reproducing an output is not the same as reproducing a conclusion. The freeze therefore also runs configured checks of the specific scientific conclusions the paper draws from the supporting bundles, for example that the Class-III label holds at every checked size or that the four map coordinates are separable. A reproducible output with a false conclusion blocks readiness. In the current freeze all 14 assets (9 paper-critical and 5 supporting) reproduce, all declared conclusion checks pass, and all seven claim groups have their evidence. The three recorded caveats are the ones stated in the paper: lens fragility of the staging classes, the absence of exponent claims, and the limited upstream bridge.

\paragraph{Provenance.} Every cached computation is keyed by its full generating configuration and by a fingerprint of the package source, configuration files, orchestration scripts, vendored Python sources, dependency versions, and the Python and NumPy versions. Within one interpreter session a cached result is not reused for a different configuration or implementation. The fingerprint is computed once per process, so the interpreter must be restarted after fingerprinted sources or configurations are changed. Unavailable observations are stored as explicit nulls and infinities as explicit tags, and neither is replaced by zero. There is one documented exception. When every shadow-ensemble member coincides (at $\lambda=1$), the Binder-type ratio is undefined but is stored as $0$, and readers should treat those entries as undefined.

\paragraph{Exact audit.} The rational certificates of Section~\ref{sec:exact-results} are produced by \path{scripts/audit_canonical_mathematics.py}. The script uses an independent implementation (\path{src/layerbirth/exact_audit.py}) that rebuilds the kernels from the configuration files and does not read any cached result. Its output, \path{notes/math_review/canonical_exact_audit.json}, contains every grid value as an exact fraction, together with configuration hashes and the maximum discrepancy from the production code. The Lean statement of Section~\ref{sec:exact-results} is in the vendored library under \path{vendor/six-birds-pica-ma_snapshot_v71/lean}. Its axiom check is \path{notes/math_review/lean_axioms.lean}.

\paragraph{Tests and figures.} The test suite (245 tests) includes regression tests for each failure mode discussed in Appendix~\ref{app:metrics-rubric}, for example unconverged stationary laws, one-way flux, fallback affinities, and incomplete control panels. The figures are generated by \path{scripts/make_paper_figures.py} from the exact audit and the frozen bundles.

\begin{table}[tbp]
\centering
\small
\caption{Claims and their primary supporting assets.}
\label{tab:appendix-claim-ledger}
\begin{tabularx}{\linewidth}{@{}L{0.33\linewidth}Y@{}}
\toprule
Claim & Primary supporting asset(s) \\
\midrule
Four-class rubric (Sec.~\ref{sec:primitives-rubric}--\ref{sec:observables-evidence}) & canonical class rubric; $P4$ metric layer; exact audit \\
Class-I vs.\ Class-II (Sec.~\ref{sec:class-i-ii}) & Class-I/II consolidation dashboard; no-fake-arrow controls; exact audit \\
Class-III / IV confirmation (Sec.~\ref{sec:class-iii-iv}) & Class-III strict confirmation; Class-III and Class-IV campaigns; exact audit; lumping proof \\
Observable map (Sec.~\ref{sec:observable-map-support}) & four-class observable map ($N=64$, recomputed) \\
Shadow ensembles (Sec.~\ref{sec:class-iii-iv}) & Class-III and Class-IV shadow panels \\
Capacity/depth (Sec.~\ref{sec:observable-map-support}) & capacity post-birth dashboard and capacity--depth summary \\
Upstream overlay (Sec.~\ref{sec:observable-map-support}) & upstream feature bridge \\
\bottomrule
\end{tabularx}
\end{table}
